\documentclass[../../../ps_q.tex]{subfiles}

\begin{document}
    真空中に電場が在り，$xy$平面上での電位分布が

    \myspaceb%
    $\displaystyle%
        \phi=a(x^2-y^2)\,(a=10\,\mathrm{V}/\mathrm{m}^2)%
    $\\
    で与えられるものとする．
    
    \begin{enumerate}
        \item%
            $xy$平面上の$x>0$かつ$y>0$の領域について，%
            $\phi=0\:\mathrm{V}$，$\phi=\pm10\:\mathrm{V}$，$\phi=\pm20\:\mathrm{V}$の%
            各々に対応する等電位線の概形を図示せよ．%
            また，同じ図上に点$(1\:\mathrm{m}, \myspacea 1\:\mathrm{m})$%
            を通る電気力線の概形を描け．
        \item%
            点$(1\:\mathrm{m}, \myspacea 0\:\mathrm{m})$での%
            電場の強さと向きを求めよ．
        \item%
            $y$軸に沿ってなめらかに動ける正の点電荷P（質量$m$，電荷$q$）に対し，%
            点$(0, \myspacea -h)\,(h>0)$で初速度（$y$軸の正の向きに大きさ$v_0$）を与える．%
            Pが$y$軸正方向の無限遠へ飛び去るための$v_0$の条件を求めよ%
            （この設問のみ文字式で答えよ）．
    \end{enumerate}
\end{document}


